%%%%%%%%%%%%%%%%%%%%%%%%%%%%%%%%%%%%%%%%%%%%%%%%%%
%%%%%%%%%%%%%%%%%%%%%%%%%%%%%%%%%%%%%%%%%%%%%%%%%%

% A primeira versão (1.0) deste modelo foi desenvolvida por Mauro Moura Severino, com a ajuda de Rubens Vasconcellos Terra Neto, para o Mestrado Profissional em Poder Legislativo (MPPL) da Câmara dos Deputados a partir do Modelo Canônico de TCC com ABNTEX2, elaborado pelo "abnTeX2 group". Cada uma das demais versões (2.0, 2.1 e 3.0) foi implementada por Mauro Moura Severino a partir da versão anterior.

%%%%%%%%%%%%%%%%%%%%%%%%%%%%%%%%%%%%%%%%%%%%%%%%%%

% Template de DISSERTAÇÃO
% Versão: v-3.0
% Data: 31/12/2024

% Esta versão viabiliza a produção de documentos com grau muito elevado de conformidade com as seguintes normas:
    % ABNT NBR 6023:2018 – Informação e documentação – Referências – Elaboração
    % ABNT NBR 6024:2012 – Informação e documentação – Numeração progressiva das seções de um documento – Apresentação
    % ABNT NBR 6027:2012 – Informação e documentação – Sumário – Apresentação
    % ABNT NBR 6028:2021 – Informação e documentação – Resumo, resenha e recensão – Apresentação
    % ABNT NBR 6034:2004 – Informação e documentação – Índice – Apresentação
    % ABNT NBR 10520:2023 – Informação e documentação – Citações em documentos – Apresentação
    % ABNT NBR 14724:2024 – Informação e documentação – Trabalhos acadêmicos – Apresentação
    % IBGE. Normas de apresentação tabular. 3. ed. Rio de Janeiro: IBGE, 1993.

%%%%%%%%%%%%%%%%%%%%%%%%%%%%%%%%%%%%%%%%%%%%%%%%%%
%%%%%%%%%%%%%%%%%%%%%%%%%%%%%%%%%%%%%%%%%%%%%%%%%%

% A EDIÇÃO DESTE ARQUIVO É DE RESPONSABILIDADE DO AUTOR, QUE DEVE EDITÁ-LO CONFORME INSTRUÇÕES A SEGUIR.

%%%%%%%%%%%%%%%%%%%%%%%%%%%%%%%%%%%%%%%%%%%%%%%%%%
%%%%%%%%%%%%%%%%%%%%%%%%%%%%%%%%%%%%%%%%%%%%%%%%%%

% ------------------------------------------------
% ------------------------------------------------
% Criação de termos do glossário
% ------------------------------------------------
% ------------------------------------------------

% ------------------------------------------------
% Para os comandos a seguir, substitua os conteúdos entre as chaves.
% ------------------------------------------------
% primeiras chaves: nome de chamada (apelido/label)
% terceiras chaves: termo como deve ser impresso
% quartas chaves: descrição do termo a ser impressa
% ------------------------------------------------

% SUGESTÃO DE USO: Não delete os comandos, copie e cole o comando que vai usar no espaço indicado a seguir, depois faça a edição.

%%%%%%%%%%%%%%%%%%%%%%%%%%%%%%%%%%%%%%%%%%%%%%%%%
% Espaço sugerido para a inclusão de termos do glossário.





%%%%%%%%%%%%%%%%%%%%%%%%%%%%%%%%%%%%%%%%%%%%%%%%%
\newglossaryentry{overleaf}{
    name={Overleaf},
    description={Plataforma \textit{online} que permite o uso da linguagem \LaTeX, com muita flexibilidade de pacotes e compiladores} % Não usar ponto final após a "description".
}

\newglossaryentry{latex}{
    name={\LaTeX},
    description={Sistema de preparação de documentos baseado em TeX} % Não usar ponto final após a "description".
}

\newglossaryentry{abntex}{
    name={abnTeX2},
    description={Modelo para formatação de trabalhos acadêmicos segundo as normas da ABNT} % Não usar ponto final após a "description".
}

\newglossaryentry{exemplo}{
    name={Exemplo},
    description={Este é um exemplo de entrada de glossário para mostrar que o termo que não é citado não é impresso} % Não usar ponto final após a "description".
}

\newglossaryentry{exemplolongo}{
    name={Exemplo longo},
    description={Este é um exemplo de entrada de glossário para mostrar como se comporta uma descrição longa. \\ \lipsum[15-17]} % Não usar ponto final após a "description".
}
